\section{Population Constants and Their Computation}\label{app:E}
e case
$r=1$, where two rounds already attain the branching coefficient. The
pointer-chasing family shows that this is not a general property of
two-stage designs.

\subsection{The Value of Adaptivity at a Given Population}\label{app:population}
For a menu $\A$ and a law $P$ on $\X\times\{-1,1\}$ write
$M_\A(P)=W_\A(P)+\Delta(P)^2$, the smallest second moment of a score that is
unbiased for $t(x,y)=yd(x)$ at every profile and label. At a full-support
law, $W_\A(P)$ is also the local asymptotic minimax risk over arbitrary
estimators (Proposition~\ref{prop:local}). For a deterministic depth-$b$ tree
$\tau$ that computes $d$, let $S_\tau(x)$ be its padded path set and
$\mathcal S_\tau$ the collection of these sets
(Appendix~\ref{app:universal}).

\begin{theorem}[Population-specific bracket, full statement of Theorem~\ref{thm:bracket}]\label{thm:bracketfull}
\textup{(a)} For every menu $\A$, law $P$ and function $h$ on
$\X\times\{-1,1\}$ with $\E_P[th]\ne0$,
\[
 M_\A(P)\ \ge\ \frac{\E_P[th]^2}{\max_{\iota\in\A}\E_P\big[\E_P[h\mid\ell_\iota(X),Y]^2\big]}.
\]
In particular $M_\A(P)\ge\E_P[t^2]/\Gamma_\A(P)$, where
$\Gamma_\A(P)=\max_\iota\E_P[\E_P[t\mid\ell_\iota(X),Y]^2]/\E_P[t^2]$ is the
largest fraction of the target's energy explained by a single interview.

\textup{(b)} If a depth-$b$ tree $\tau$ computes $d$, then
$M_{\rm branch}(P)=\E_P[t^2]$ and
\[
 \frac1{\Gamma_{\rm fixed}(P)}
 \ \le\ \frac{M_{\rm fixed}(P)}{M_{\rm branch}(P)}
 \ \le\ N_P(\tau):=\frac{\Big(\sum_{S\in\mathcal S_\tau}\E_P\big[t^2\mathbf 1\{S_\tau(X)=S\}\big]^{1/2}\Big)^2}{\E_P[t^2]}
 \ \le\ |\mathcal S_\tau|\le K_{m,b}.
\]

\textup{(c)} Both inequalities in (b) are equalities, with value $D$, for the
address contrast with $\kappa=b-1$ under the uniform law with independent
uniform labels.
\end{theorem}

The quantity $N_P(\tau)=\exp H_{1/2}(w)$ is the exponential of the R\'enyi
entropy of order $1/2$ of the weights
$w_S=\E_P[t^21\{S_\tau=S\}]/\E_P[t^2]$: it is an effective number of
question sets that the tree actually uses at $P$. Adaptivity can be valuable
only when the tree's questions diverge across many well-populated paths. It
is valuable when, in addition, no fixed set of $b$ questions explains most of
the target, so that $\Gamma_{\rm fixed}(P)$ is small; correlated features and
shared coordinates raise $\Gamma_{\rm fixed}$. The variance ratio is
recovered as $\rho_P=(M_{\rm fixed}-\Delta^2)/(\E_Pt^2-\Delta^2)$.

\begin{proof}
(a) Let $H$ be any feasible score with interview probabilities $q_\iota$ and
transcript tables $g_\iota$, so $\sum_\iota g_\iota(\ell_\iota(x),y)=t(x,y)$
everywhere and $\E[H^2]=\sum_\iota\|g_\iota\|_P^2/q_\iota\ge(\sum_\iota\|g_\iota\|_P)^2$
by Cauchy--Schwarz in $q$. For every $h$,
\[
 \E_P[th]=\sum_\iota\E_P\big[g_\iota\,\E_P[h\mid\ell_\iota(X),Y]\big]
 \le\Big(\sum_\iota\|g_\iota\|_P\Big)\max_\iota\big\|\E_P[h\mid\ell_\iota(X),Y]\big\|_P,
\]
because $g_\iota$ is a function of $(\ell_\iota(X),Y)$. Squaring and taking
the infimum over feasible scores proves the display; $h=t$ gives the second
statement.

(b) Following $\tau$ returns $t$ itself, and Jensen's inequality gives
$\E[H^2]\ge\E[\E[H\mid X,Y]^2]=\E_P[t^2]$ for every unbiased score. The lower
bound is (a) with $\A$ the fixed forms. For the upper bound choose
$q_S>0$ on $\mathcal S_\tau$, draw $S$ before the person answers, and return
$t(X,Y)/q_S$ if the tree's path completes inside $S$ and zero otherwise.
Completion occurs exactly when $S=S_\tau(X)$
(Appendix~\ref{app:universal}), so the score is unbiased at every profile and
label, and its second moment is
$\sum_Sq_S^{-1}\E_P[t^2\mathbf 1\{S_\tau(X)=S\}]$. Minimizing over the simplex
gives $q_S\propto\E_P[t^21\{S_\tau=S\}]^{1/2}$ and the value
$N_P(\tau)\E_P[t^2]$. Sets of zero weight are handled by letting their
probability tend to zero. The last inequalities are the power-mean
inequality and $|\mathcal S_\tau|\le K_{m,b}$.

(c) Under the uniform law the weights are equal over the $D$ path sets, so
$N_P=D$. With independent labels,
$\E[t\mid X_T,Y]=Y\E[d\mid X_T]$, whose energy is at most $2^{-\kappa}=1/D$
by~\eqref{eq:addressinformation}. Hence $1/\Gamma_{\rm fixed}\ge D$, and the
universal bound $M_{\rm fixed}\le D$ forces equality throughout.
\end{proof}

\paragraph{The bracket on the variance scale.}
Section~\ref{sec:experiments} reports the number of records a design needs
relative to branching, that is the ratio of variances
$W_\A(P)=M_\A(P)-\Delta(P)^2$ of pointwise-unbiased designs rather than the
ratio of second moments. The two scales differ when $\Delta(P)\neq0$, and the
bracket transfers as follows.

\begin{corollary}[Bracket on the variance scale]\label{cor:bracketvar}
Under the assumptions of Theorem~\ref{thm:bracketfull}(b), with
$\sigma_P^2=\E_P[t^2]-\Delta(P)^2=W_{\rm branch}(P)>0$,
\[
 \frac{\E_P[t^2]/\Gamma_{\rm fixed}(P)-\Delta(P)^2}{\sigma_P^2}
 \ \le\ \frac{W_{\rm fixed}(P)}{W_{\rm branch}(P)}
 \ \le\ \frac{N_P(\tau)\,\E_P[t^2]-\Delta(P)^2}{\sigma_P^2}.
\]
When the second-moment endpoints exceed one, the variance-scale endpoints are
larger than them, so the variance ratio is bracketed by a wider interval
than the second-moment ratio.
\end{corollary}
\begin{proof}
Every design in \eqref{eq:mainW} is unbiased for $\Delta(P)$, so its variance
is its second moment minus $\Delta(P)^2$; in particular
$W_{\rm fixed}=M_{\rm fixed}-\Delta^2$ and $W_{\rm branch}=\E_P[t^2]-\Delta^2$.
Subtract $\Delta^2$ from the three terms of Theorem~\ref{thm:bracketfull}(b)
multiplied by $\E_P[t^2]$ and divide by $\sigma_P^2$. The last statement is
the inequality $(a-c)/(b-c)\ge a/b$ for $a\ge b>c\ge0$.
\end{proof}
The tables of Section~\ref{sec:experiments} and Appendix~\ref{app:screening}
report this variance-scale bracket, and "inside the bracket" refers to it.


\paragraph{Contrasts that branching cannot compute.}
For an ensemble whose contrast is not computed by one depth-$b$ tree, part (a)
applies to both menus. Every depth-$b$ branching interview is a coarsening of
a complete depth-$b$ tree, and the energy of a tree is a sum over its leaves.
Hence $\max_\iota\E_P[\E_P[h\mid\ell_\iota,Y]^2]$ over all depth-$b$ trees is
computed exactly by dynamic programming over partial assignments, in time
$O(m^b2^b)$. This yields certified lower bounds on $M_{\rm branch}(P)$ for every
dual witness $h$. Appendix~\ref{app:forestexact} uses these bounds inside
column generation to compute the branching coefficient over all depth-three
trees.

\subsection{A Computable Upper End: the Optimized Path-Set Design}\label{app:upper}
The upper end $N_P(\tau)$ of Theorem~\ref{thm:bracket} is the second moment of completion-only path guessing, which discards every answer on a form except a completed path. Optimizing the full score program on the same small menu tightens it without enumerating the $\binom mb$ forms.

\begin{proposition}[Optimized path-set design]\label{prop:upper}
Let a depth-$b$ tree $\tau$ compute $d$, let $S_1,\ldots,S_L$ be its distinct padded path sets ($L\le K_{m,b}$), and let $M_{\rm PS}(P)$ be the optimum of the population program \eqref{eq:mainW} over the fixed-form menu $\{S_1,\ldots,S_L\}$, with all pointwise-unbiased tables. Then
\[
 M_{\rm fixed}(P)\ \le\ M_{\rm PS}(P)\ \le\ N_P(\tau)\,M_{\rm branch}(P),
\]
and the same chain holds on the variance scale after subtracting $\Delta(P)^2$ and dividing by $\sigma_P^2$.
\end{proposition}
\begin{proof}
The menu $\{S_1,\ldots,S_L\}$ is a set of fixed forms, so its optimum is at least the optimum over all forms. The menu also identifies $d$: every leaf indicator of $\tau$ depends only on that leaf's path coordinates, which lie in its padded set, so $d$ is a sum of tables on the $S_j$. Completion-only path guessing with $q_j\propto w_j^{1/2}$ is a feasible design in this program, and its second moment is $(\sum_jw_j^{1/2})^2\E_P[t^2]=N_P(\tau)\E_P[t^2]$ by \eqref{eq:pathvariance}; Theorem~\ref{thm:bracket} gives $M_{\rm branch}(P)=\E_P[t^2]$. Every design in the program is pointwise unbiased, so all three quantities exceed their variances by the same $\Delta(P)^2$.
\end{proof}
The program has at most $2L2^b$ score variables and is solved from reference marginals; Appendix~\ref{app:screendetail} compares $M_{\rm PS}$ with the all-forms optimum on every screening run. The two coincide whenever an optimal fixed design uses only path sets of $\tau$; in general the gap measures the value of forms that no path of the tree reads.

\subsection{Computing Optimal Designs for Every Protocol Class}\label{app:engine}
All population constants in Section~\ref{sec:experiments} are computed by one column-generation solver whose cost depends on $\binom mb2^b$, not on the $2^m$ profiles.

\paragraph{Cells and masses.}
Pad every transcript of an interview to exactly $b$ coordinates. An interview of $\A_r$ is then a partition of the cube into cells $C(T,u)=\{x:x_T=u\}$ with $|T|=b$; for a fixed form all cells share $T$, for an $r$-round interview $T$ varies across cells. Padding refines a transcript and therefore never increases the optimal second moment. For a law $P$ that mixes records with a uniform mass $\alpha$ over profiles and labels, the masses $P(T,u,z)$ of all $\binom mb2^{b+1}$ cells are computed once from the records.

\paragraph{Master problem.}
For a list of interviews, minimize $\sum_{\iota,c,z}P(c,z)g_{\iota,z}(c)^2/q_\iota$ subject to $\sum_\iota q_\iota=1$ and the Fourier identities $\sum_{\iota,c}\widehat{\mathbf 1}_c(S)g_{\iota,z}(c)=(2z-1)\widehat d(S)$ for all $|S|\le b$, where $\widehat{\mathbf 1}_{C(T,u)}(S)=2^{-b}\chi_S(u)$ for $S\subseteq T$ and zero otherwise. Because cell indicators and $d$ have degree at most $b$, the identities are equivalent to pointwise unbiasedness at every profile. The problem is a second-order cone program; we solve it with Clarabel through CVXPY at tolerance $10^{-10}$. Inside column generation the solver\textquotesingle s solution steers the iteration; the final design is repaired to exact feasibility by the least-norm correction weighted by the interview probabilities, which adds the least second moment, so every reported upper bound is the second moment of an exactly pointwise-unbiased design. Looser solver tolerances produced visibly suboptimal masters in our tests, which the consistency checks of Appendix~\ref{app:F} (the fixed optimum may not exceed the path-set or fitted designs) detect.

\paragraph{Pricing.}
Let $\eta_z(S)$ be the multipliers of the identities. The information of a cell is
\[
 e(T,u)=\sum_z\frac{\psi_z(T,u)^2}{P(T,u,z)},\qquad \psi_z(T,u)=2^{-b}\sum_{S\subseteq T}\eta_z(S)\chi_S(u),
\]
a Walsh transform of $\eta$ restricted to $T$; all cells cost $O(\binom mbb2^b)$ operations. The information of an interview is the sum over its cells. The most informative interview in $\A_r$ is found by dynamic programming over partial assignments $(A,a)$, $|A|<b$:
\[
 V_1(A,a)=\max_{T\supseteq A,|T|=b}\ \sum_{u:\,u_A=a}e(T,u),\qquad
 V_\rho(A,a)=\max_{B\supsetneq A,\,|B|\le b}\ \sum_{a':\,a'_A=a}V_{\rho-1}(B,a'),
\]
with $V_\rho(T,u)=e(T,u)$ for $|T|=b$; the best $r$-round interview has information $V_r(\emptyset,\emptyset)$, the best tree is the special case of singleton batches, and the best form is $\max_T\sum_ue(T,u)$. The program visits $O(\sum_{k\le b}\binom mk4^k)$ entries, and an optimal interview is reconstructed by following the maximizers.

\paragraph{Certificates and stabilization.}
Normalize $\eta$ so that $\sum_{z,S}\eta_z(S)(2z-1)\widehat d(S)=1$. By the last display of Theorem~\ref{thm:bracket}, $1/\max_{\iota\in\A_r}e(\iota)$ is a lower bound on $M_{\A_r}(P)$ for every such $\eta$, and the master value is an upper bound. Every pointwise-unbiased design also satisfies $M_\A(P)\ge\E_P[t^2]$ by Jensen\textquotesingle s inequality, which certifies designs that attain full information. Each iteration adds up to six (twelve for the forests and the 21-item study) interviews whose information exceeds the master\textquotesingle s maximum, and the solver stops when no interview of the class improves the master or when the bounds agree to a relative error of $10^{-7}$ ($10^{-4}$ for the forest tasks at $b\ge4$). Master duals of column generation are degenerate, and unstabilized pricing can stall for many iterations; we price at a convex combination, with equal weights, of the current dual and the dual with the best dual bound so far \citep{wentges1997}, and fall back to the current dual when the smoothed point finds no improving interview. Removing unused columns from the master made the iteration cycle in our tests, so the master keeps every generated interview. The master starts from the support of the optimal fixed design (and, for the forests, the padded component trees), so it stays small.

\paragraph{Validation.}
On stored tasks the solver reproduces the fixed constants, $\Gamma_P$ and the certified branching constants of an independent dense implementation (released with the code) to six significant digits, and every generated interview is checked to partition the cube. Branching and $r$-round designs at $m\le12$ converge in tens to a few hundred iterations; pricing is no longer the bottleneck, and the master problem dominates the running time (Appendix~\ref{app:forestexact}).

\subsection{The Form Program and the Path-Set Menu}\label{sec:pathdesign}
A form can reveal useful coordinates even when it fails to complete a tree path. For a frozen depth-$b$ tree, take the coordinate set on each root-to-leaf path, extend it to size $b$ with fixed unused coordinates, and retain the $L$ distinct sets as a tractable menu. Let $g_{T,y}(u)$ be a table on \emph{all} answers $u$ to form $T$, and let $\widehat P_T(u,y)$ be calibrated marginal frequencies. For binary features solve
\begin{align}
 \min_{q,g}\quad&\sum_{T,u,y}\widehat P_T(u,y)
                       \frac{g_{T,y}(u)^2}{q_T},\label{eq:fullprogram}\\
 \text{subject to}\quad&\sum_{T\supseteq S}\widehat g_{T,y}(S)
                   =y\widehat d(S)\quad\text{for every }S,y,\nonumber\\
 &\sum_Tq_T=1,\qquad q_T\ge\beta/L.\nonumber
\end{align}
Here $\widehat g_{T,y}(S)$ is the Fourier coefficient of the table on its own form; it is a linear transform of $g$. Fitted designs in Section~\ref{sec:experiments} use $\beta=0.05$. These constraints are exactly equivalent to $\sum_Tg_{T,y}(x_T)=yd(x)$. Returning $g_{T,Y}(X_T)/q_T$ is therefore globally unbiased even when the reference frequencies are inaccurate; the reference affects only the variance.

There are at most $2L2^b$ score entries and coefficient constraints, and the program avoids enumerating the $2^m$ full profiles. Pairwise shared-coordinate corrections span the complete unbiased score space on the supplied menu over the full binary cube (Proposition~\ref{prop:sharing}). A correction adds $f(X_{T\cap T'},Y)/q_T$ on form $T$ and subtracts it divided by $q_{T'}$ on form $T'$. The optimization therefore includes every such control variate. In the language of coarsening at random, \eqref{eq:fullprogram} chooses the optimal augmentation of the inverse-probability-weighted score under the reference law, jointly with the form probabilities. For ordinal data the same program applies to tensor contrasts on the categories defined by the frozen tree's thresholds.

For an existing randomized archive, keep its known $q$ fixed and optimize only the tables. Form-specific marginal frequencies can be estimated from incomplete training records; they need not be assembled into a joint reference distribution. Identification then requires coverage by the actual form menu, not only the degree bound for the menu of all size-$b$ forms.

A \emph{completion-only} score returns an inverse-probability weighted target on a matching padded path and zero otherwise; the experiments call its uniform version path guessing.

\paragraph{Proper scoring rules.}
For Brier improvement, with $z=(Y+1)/2$,
\[
 (p_0-z)^2-(p_1-z)^2=d(1-p_0-p_1)+Yd.
\]
In the form program, replace $Yd$ by this known target and impose separate coefficient identities for each label. The same unbiasedness guarantee holds. For ensembles of degree $r$, Brier loss can require degree $2r$; log loss can introduce still higher interactions. Appendix~\ref{app:metrics} gives the identification conditions and explains why AUC requires a separate pairwise design.


\subsection{Shared-Coordinate Corrections and the Complete Score Space}\label{app:fullscores}
Fix a menu of coordinate sets $\mathcal T$. For binary features write
$g_{T,y}(x_T)=\sum_{S\subseteq T}\gamma_{T,y,S}\chi_S(x)$.
Uniqueness of the Fourier expansion gives the following elementary
representation.

\begin{proposition}[All shared-coordinate corrections]\label{prop:sharing}
The score tables satisfy $\sum_Tg_{T,y}(x_T)=yd(x)$ on the full binary cube
if and only if
\[
       \sum_{T\supseteq S}\gamma_{T,y,S}=y\widehat d(S)
       \quad\text{for every }S.
\]
For a fixed positive design $q$, every change to a globally unbiased score
is a sum of pairwise transfers of functions on shared coordinates:
add $f(X_{T\cap T'},Y)/q_T$ to form $T$ and subtract the same function
divided by $q_{T'}$ from form $T'$.
\end{proposition}
\begin{proof}
Equate the coefficients of each $\chi_S$. For the second claim, any score
difference has, for each $S$ and each label, coefficients summing to zero
over forms containing $S$. Choose one of those forms as an anchor and
transfer every other coefficient to or from the anchor. Each transfer is
observable under both forms. Terms appearing in only one form must have
zero difference. Summing the transfers proves the statement.
\end{proof}
Thus including all such terms uses the entire unbiased score space on the
supplied menu. This is a product-basis implementation of standard control
variates and known-coarsening influence-function optimization, not a new
general missing-data efficiency bound.

The implementation works in cell values. A block of the constraint matrix
is the normalized Walsh transform of a form's table, placed in the rows of
its contained coordinate subsets. If $F$ is the resulting matrix, the
constraints are $Fg_y=y\widehat d$. Its rows are orthogonal:
$FF^{\mathsf T}$ is diagonal, with entry
$\sum_{T\supseteq S}2^{-|T|}$ for subset $S$. We use this identity to
repair numerical equality residuals explicitly. Rotated second-order cones
represent $g_{T,y}(u)^2\le q_Tv_{T,y}(u)$.

For a given population, marginal form--label frequencies are sufficient for
the objective. No full-profile enumeration is needed. With $L$ forms of
size at most $b$, there are at most $2L2^b$ score entries and $2L2^b$
coefficient constraints. A tree provides the target coefficients through
its leaves, or by Fourier transforming its known completion-score
decomposition. Choosing an unrestricted menu can still require
$\binom mb$ forms; this implementation does not solve that combinatorial
design problem.

For ordinal features one splits each coordinate at all thresholds used for it
in the frozen tree, with a fixed middle cut for otherwise unused padding
coordinates. We use orthonormal Helmert contrasts on those categories and
their tensor products. The same argument is exact on the resulting product
space. This retains every observed tree-threshold category, but not
within-category distinctions between raw ordinal answers. Its cell count
is $\sum_T\prod_{j\in T}s_j$, where $s_j$ is the category count.


\subsection{Efficiency and Augmented Fixed Forms}\label{app:augmentation}
\subsubsection{Relation to the Missing-Data Efficiency Bound}
Fix a known form law $q$ and a full-support law $P$ on the finite complete-record space. The observed record is $O=(A,o_A(Z))$. Its score along a centered complete-data perturbation $h$ is $\E_P[h(Z)\mid O]$. A mean-zero observed influence function $\phi(O)$ for $\Delta=\E_Pt$ must therefore satisfy
\[
 \E[\phi(O)\mid Z]=t(Z)-\Delta.
\]
Conversely this equation gives the required pairing with every complete-data score, and hence an influence function in this finite regular model. Minimizing $\E\phi^2$ is the efficiency bound. Substituting $\phi=H-\Delta$ identifies it with the fixed-design population variance $W_q(P)$. This is the usual known-coarsening influence-function geometry \citep{robins1994,tsiatis2006}, not a new asymptotic efficiency formula. The proposition's additional statement is finite-$n$ pointwise optimality over globally unbiased nonlinear estimators and sequential form assignments, together with its extension to boundary laws. It does not assert a finite-sample bound for all locally regular biased estimators.

Fixed random forms implement planned missingness. A tree's terminal cell is generated using observed answers, so its known acquisition mechanism is a designed coarsening-at-random mechanism \citep{heitjan1991,gill1997}. This description does not cover arbitrary unrecorded selection or nonresponse.

\subsection{Learning a Design from a Compact Tree}\label{app:pathdesign}
Fix a tree that recovers $t(z)=yd(x)$ within $b$ distinct queries. Pad each path deterministically to $b$ coordinates. Enumerate its distinct path sets $\mathcal S=\{S_1,\ldots,S_L\}$ and let $S(z)$ denote the reached set. If an assigned fixed form $S_j$ permits following the padded tree to completion, then $S(z)=S_j$, as in Appendix~\ref{app:universal}. For any strictly positive form law $q$, the score
\[
 H_q(z,J)=\frac{t(z)}{q_J}\mathbf1\{S(z)=S_J\},\qquad J\sim q,
\]
is computable from the selected form and label. It is unbiased conditional on every full record. With $w_j=\E_P[t(Z)^2\mathbf1\{S(Z)=S_j\}]$,
\begin{equation}
 \Var_P H_q=\sum_{j=1}^L\frac{w_j}{q_j}-\Delta(P)^2.
 \label{eq:pathvariance}
\end{equation}
Within this importance-score family the unconstrained optimum is $q_j\propto\sqrt{w_j}$, with infimal second moment $(\sum_j\sqrt{w_j})^2$. Zero-mass cells are handled by limits. We enforce $q_j\ge\beta/L$, ensuring global unbiasedness even on path sets absent from calibration. If some $w_j>0$, the constrained optimum satisfies $q_j=\max\{\beta/L,c\sqrt{w_j}\}$, with $c$ chosen to normalize the probabilities. If all weights vanish, every feasible $q$ is optimal. Iteratively fixing cells below the floor computes the positive-weight solution without profile enumeration. This oracle optimizes a specified score family, not all fixed-form estimators.

\subsubsection{A Worked Two-Question Example}
Let $d(x)=[(1+x_1)x_2+(1-x_1)x_3]/2$. Branching queries $x_1$ and then $x_2$ when $x_1=1$, or $x_3$ when $x_1=-1$. The two fixed forms are
\begin{center}
\begin{tabular}{lll}
\toprule
Form & Probability & Observed score\\
\midrule
$\{1,2\}$ & $1/2$ & $2Yx_2\mathbf1\{x_1=1\}$\\
$\{1,3\}$ & $1/2$ & $2Yx_3\mathbf1\{x_1=-1\}$\\
\bottomrule
\end{tabular}
\end{center}
At every full record the mean score is $Yd$ and its second moment is two. The branching score has second moment one. For an independent fair label both score means vanish, so their sample-mean variances are exactly $2/n$ and $1/n$. This describes these concrete estimators; the minimax results additionally justify their worst-case comparison against all protocols.

\subsubsection{A Finite Calibration Guarantee}
Suppose the tree is fitted independently of $r$ calibration records and $|t|\le M$. Full calibration records reveal $W_j=t^2\mathbf1\{S=S_j\}$ directly. Alternatively, assign each calibration person a uniformly random fixed path set. On successful completion return $W_j=L t^2\mathbf1\{J=j,S=S_j\}$ and return zero on failure. This alternative obeys the same fixed per-person cap and uses only observable quantities. In both cases $\E W=w$. Let $\widehat w$ be the mean and let $\widehat q$ minimize $\sum_j\widehat w_j/q_j$ subject to $q_j\ge\beta/L$.

\begin{proposition}[Learning path-set probabilities]\label{prop:pilot}
Let $q_\beta^\star$ be the population optimum with the same probability floor. Set $B=M^2$ for full calibration or $B=LM^2$ for uniformly assigned fixed calibration. With probability at least $1-\alpha$ over the $r$ independent calibration records,
\[
 \Var_P H_{\widehat q}-\Var_P H_{q_\beta^\star}
 \le\frac{2LB}{\beta}\sqrt{\frac{2(L\log2+\log(1/\alpha))}{r}}.
\]
The resulting audit mean on fresh respondents remains globally unbiased for every realized calibration sample.
\end{proposition}
\begin{proof}
Each calibration vector has at most one nonzero entry and $\|W\|_1\le B$. For any sign vector $\sigma\in\{-1,1\}^L$, $\sigma^TW\in[-B,B]$. Hoeffding's inequality and a union bound over $2^L$ vectors give
\[
 \|\widehat w-w\|_1
 =\max_\sigma\sigma^T(\widehat w-w)
 \le B\sqrt{2(L\log2+\log(1/\alpha))/r}.
\]
The difference between the empirical and population second-moment objective is uniformly at most $(L/\beta)\|\widehat w-w\|_1$. Empirical minimization costs at most twice this uniform deviation. Subtracting the common $\Delta^2$ proves the bound. Pointwise conditional unbiasedness holds for each positive $q$, hence also conditional on the entire calibration sample.
\end{proof}
The bound is conservative and becomes informative only with sufficiently large calibration samples. It proves consistency relative to a constrained family oracle, not a calibration-free optimality claim. Form enumeration takes time polynomial in the given tree's size and $m$; that tree may itself have exponentially many leaves in $b$. A deployment needs only its form probabilities and the tree, not a table over all $s^m$ profiles.

\subsubsection{Confidence Intervals and Costs}
Conditional on calibration, each path score has magnitude at most $M/q_{\min}$. Thus the ordinary score mean has a distribution-free $(1-\alpha)$ confidence interval with half-width
\[
 \frac{M}{q_{\min}}\sqrt{\frac{2\log(2/\alpha)}n},
\]
intersected with $[-M,M]$. The branching score has the same bound with $q_{\min}=1$. These Hoeffding intervals are valid but conservative; they are not claims of optimal interval width or power. The paper's risk comparisons concern MSE and do not order intervals for biased estimators.

For a menu with known per-person costs $c_\iota>0$, let $c(q)=\sum_\iota q_\iota c_\iota$. An iid unbiased score mean with $n$ respondents uses expected cost $nc(q)$, so its worst-case MSE-times-cost upper criterion is $c(q)\max_x\sum_\iota g_\iota(x)^2/q_\iota$. Set $r_\iota=q_\iota/c(q)$. Joint cost-weighted design becomes the convex program
\[
 U_c(d,\A)=\inf_{\substack{r,g:\,\sum_\iota c_\iota r_\iota=1\\ \sum_\iota g_\iota=d}}
              \max_x\sum_\iota\frac{g_\iota(x)^2}{r_\iota}.
\]
Conversely $q_\iota=r_\iota/\sum_jr_j$ recovers the design, proving equivalence of the two criteria. For fixed forms, $c_\iota$ may include recruitment, the label and item-specific costs. For a branching tree, taking $c_\iota$ as its maximum path cost gives a conservative criterion. This uses expected total expenditure and continuous sample-size scaling; it does not impose an almost-sure stopping budget or establish a sharp adaptivity ratio for arbitrary costs.

\subsection{Learning the Design from the Audit Itself}\label{app:selfcal}
The population-specific constants $W_\A(P)$ are attained by designs tuned to
$P$. Section~\ref{sec:experiments} tunes them on separate reference records.
No reference is needed asymptotically: because every cell mass of every
interview is estimable from the records assigned that interview, a pilot
phase of the audit itself can supply the design, and pointwise unbiasedness
keeps every intermediate design exactly unbiased.

For a finite menu $\A$, a law $P$ on $\X\times\{-1,1\}$ and a design $(q,g)$
feasible for \eqref{eq:mainW}, write $J_P(q,g)=\sum_\iota\sum_{c,y}P_\iota(c,y)g_\iota(c,y)^2/q_\iota$
for its second moment, where $P_\iota(c,y)=P(\ell_\iota(X)=c,Y=y)$ are the
cell masses, so that $M_\A(P)=\inf_{q,g}J_P(q,g)$. For $\beta\in(0,1]$ let
$M_\A^{\beta}(P)$ denote the infimum over designs with $q_\iota\ge\beta/|\A|$
for all $\iota$.

\begin{proposition}[Two-phase self-calibration]\label{prop:selfcal}
Let $n_1=n_1(n)\to\infty$ with $n_1/n\to0$, and let $\alpha_n,\beta_n\in(0,1]$
tend to zero with $\alpha_n n_1^{1/2}\to\infty$. Phase one interviews $n_1$
records with the uniform design on $\A$ and forms, for every interview $\iota$,
the Horvitz--Thompson estimates
$\widehat P_\iota(c,y)=(|\A|/n_1)\sum_{i\le n_1}\mathbf 1\{\iota_i=\iota,\ \ell_\iota(X_i)=c,\ Y_i=y\}$
and the smoothed masses $\widetilde P_\iota=(1-\alpha_n)\widehat P_\iota+\alpha_n u_\iota$,
where $u_\iota$ is uniform on the cells and labels of $\iota$. Let
$(\hat q,\hat g)$ minimize $J_{\widetilde P}$ over feasible designs with
$q\ge\beta_n/|\A|$, and interview the remaining $n_2=n-n_1$ records with
$(\hat q,\hat g)$. Then the mean $\widehat\Delta$ of their scores is
unbiased for $\Delta(P)$ conditionally on phase one, for every $P$, and
\[
 n_2\Var(\widehat\Delta\mid\text{phase one})=J_P(\hat q,\hat g)-\Delta(P)^2,
 \qquad J_P(\hat q,\hat g)\ \xrightarrow{\ \PP\ }\ M_\A(P),
 \qquad n\,\Var(\widehat\Delta\mid\text{phase one})\xrightarrow{\ \PP\ }W_\A(P).
\]
\end{proposition}
\begin{proof}
Unbiasedness holds because $(\hat q,\hat g)$ satisfies the pointwise
constraints $\sum_\iota\hat g_\iota(\ell_\iota(x),y)=t(x,y)$ for every $(x,y)$,
whatever the phase-one data. Conditionally on phase one the phase-two scores
are iid with mean $\Delta(P)$ and second moment $J_P(\hat q,\hat g)$, which
gives the variance identity. For the limit put $P^{\alpha}_\iota=(1-\alpha_n)P_\iota+\alpha_nu_\iota$.
Each $\widehat P_\iota(c,y)$ is an average of $n_1$ iid bounded terms with
mean $P_\iota(c,y)$, so $\max_{\iota,c,y}|\widehat P_\iota(c,y)-P_\iota(c,y)|=O_\PP(n_1^{-1/2})$
over the finitely many cells. Since
$\widetilde P_\iota-P^\alpha_\iota=(1-\alpha_n)(\widehat P_\iota-P_\iota)$ and
$P^\alpha_\iota\ge\alpha_nu_\iota$,
\[
 \rho_n:=\max_{\iota,c,y}\Big|\frac{\widetilde P_\iota(c,y)}{P^\alpha_\iota(c,y)}-1\Big|
 \le\frac{\max_{\iota,c,y}|\widehat P_\iota(c,y)-P_\iota(c,y)|}{\alpha_n\min_\iota u_\iota}
 =O_\PP\big((\alpha_nn_1^{1/2})^{-1}\big)\to0 .
\]
Let $(q^\alpha,g^\alpha)$ be optimal for $J_{P^\alpha}$ under the floor
$\beta_n$. Because $J$ is linear in the masses,
\[
 (1-\alpha_n)J_P(\hat q,\hat g)\le J_{P^\alpha}(\hat q,\hat g)\le\frac{J_{\widetilde P}(\hat q,\hat g)}{1-\rho_n}
 \le\frac{J_{\widetilde P}(q^\alpha,g^\alpha)}{1-\rho_n}\le\frac{1+\rho_n}{1-\rho_n}\,M^{\beta_n}_\A(P^{\alpha}) .
\]
Finally $M^{\beta}_\A(P^{\alpha})\to M_\A(P)$ as $\alpha,\beta\to0$: for any
feasible $(q,g)$, mixing $q$ with the uniform design at weight $\beta$ keeps
the tables feasible and multiplies $J$ by at most $(1-\beta)^{-1}$, so
$M_\A^\beta(P^\alpha)\le(1-\beta)^{-1}J_{P^\alpha}(q,g)\le(1-\beta)^{-1}[(1-\alpha)J_P(q,g)+\alpha J_{u}(q,g)]$,
whose infimum over $(q,g)$ tends to $M_\A(P)$; conversely
$M_\A^\beta(P^\alpha)\ge(1-\alpha)M_\A(P)$. Combining the displays gives
$\limsup J_P(\hat q,\hat g)\le M_\A(P)$ in probability, and
$J_P(\hat q,\hat g)\ge M_\A(P)$ holds for every feasible design.
Since $n_2/n\to1$, the last claim follows from
$W_\A(P)=M_\A(P)-\Delta(P)^2$.
\end{proof}

The proposition is an asymptotic statement about a finite menu: phase one
must estimate the masses of all $\sum_\iota|{\rm cells}(\iota)|$ cells, which
requires $n_1$ large compared with $|\A|2^b$, and in the experiments of
Section~\ref{sec:experiments} reference records play the role of phase one.
A sequential version that refits the design after every batch of records
inherits exact unbiasedness by the same argument; its rate is not analyzed
here. Learning the design across records is standard adaptive importance
sampling; the point here is that the pointwise constraints make every fitted
design exactly unbiased, so that the price of learning is paid in variance
only.

